Welcome to The Fourth Arabic Natural Language Processing Conference (ArabicNLP 2026), co-located with EMNLP 2026 in Budapest, Hungary. ArabicNLP 2026 is the eleventh edition of the WANLP/ArabicNLP meeting series, which has developed a growing reputation as a high quality venue for researchers and engineers working on Arabic NLP, where they share and discuss their ongoing work. The first in the WANLP series was held in Doha, Qatar (EMNLP 2014), followed by Beijing, China (ACL 2015), Valencia, Spain (EACL 2017), Florence, Italy (ACL 2019), online with COLING 2020, online with EACL 2021, a hybrid event in Abu Dhabi, UAE (EMNLP 2022), and then in-person events in Singapore with EMNLP 2023, Bangkok, Thailand with ACL 2024, and Suzhou, China with EMNLP 2025.

For this year's edition of ArabicNLP, we received a record 193 valid main conference submissions, more than double the 95 received last year. Of these, 164 received at least one review, the remainder having been desk rejected before review for anonymity, formatting or completeness violations, and 139 went through the full review process, in which each paper was read by at least three reviewers and then by one Area Chair, a role newly added to ArabicNLP this year. We accepted 37 papers, for an overall acceptance rate of 22.56\% across all submissions that received at least one review. This is the most selective edition to date in the WANLP/ArabicNLP series.

The accepted papers are presented over the two days of the conference in five oral sessions and two poster sessions: 21 papers as oral presentations and 16 as posters. The oral sessions are organized around the themes that emerged most strongly from the accepted set: Arabic speech, resources and benchmarks; dialects, machine translation and content moderation; Arabic for education and Islamic knowledge; Arabic resources, morphology and syntax; and LLM evaluation, safety and detoxification. Alongside the main conference papers, ArabicNLP 2026 hosts fifteen shared tasks with 212 papers (15 overview papers and 197 system description papers), which appear in Volume 2 of these proceedings.

ArabicNLP 2026 includes an invited talk by Mona Diab (Carnegie Mellon University), entitled "From NLP for Arabic to NLP Through Arabic: What Must AI Understand Before We Trust It to Act?", a fireside chat, and two panel discussions. We also recognize outstanding contributions through a Best Paper Award and a Best Resource Award, announced during the conference.

We were able to secure sponsorship funding from different institutions: Google Research, QatarDebate Center (Qatar), Aram Lab (MBZUAI), CAMeL Lab (NYU Abu Dhabi), Tuwaiq Academy (Saudi Arabia), and Resemble AI. We used the sponsorship funds to support student registrations. We thank all our sponsors for their generous support and their help in building up the Arabic NLP community. Finally, we extend our gratitude to everyone who submitted a paper to the conference, and to the Area Chairs and Program Committee members for their diligent efforts in providing reviews within a very tight time frame.

Muhammad Abdul-Mageed, General Chair, on behalf of the conference organizers.

Website of the conference: https://arabicnlp2026.sigarab.org
